%! AP Statistics — Unit 5, Lesson 5.1 — Homework (Answer Key)
\documentclass[10pt]{article}
\usepackage{apstats-article}
\usepackage{apstats-key}   % pulls in -boxes; do NOT also load -boxes
\newcommand{\TallMath}[1]{$\displaystyle #1\rule[-1.4em]{0pt}{3.2em}$}

\begin{document}

\pageheader{Unit 5, Lesson 5.1}{Homework --- Answer Key}
\namedateperiod

\vspace{0.15in}

\begin{enumerate}

  \item
        \begin{enumerate}[label=\alph*.]
          \item
                \begin{itemize}
                  \item Statistic: \ans{sample mean heart rate, 72 bpm} \quad Notation: \ans{$\bar{x}$}
                  \item Parameter: \ans{true mean heart rate of all adults} \quad Notation: \ans{$\mu$}
                \end{itemize}
          \item
                \begin{itemize}
                  \item Statistic: \ans{sample proportion who exercise $\ge$ 3 times/week, $64/200 = 0.32$} \quad Notation: \ans{$\hat{p}$}
                  \item Parameter: \ans{true proportion of all college students who exercise $\ge$ 3 times/week} \quad Notation: \ans{$p$}
                \end{itemize}
        \end{enumerate}

  \item
        \begin{enumerate}[label=\alph*.]
          \item \ansline{No --- they are two different sample means from two different randomly selected groups of seedlings, so they are not the same measurement. Each describes only the 15 seedlings in that sample.}
          \item \ansline{No. Different random samples from the same greenhouse will naturally produce different sample means by chance (sampling variability). Neither measurement is wrong; they reflect natural variation.}
          \item \ans{sampling variability}
        \end{enumerate}

  \item
        \begin{enumerate}[label=\alph*.]
          \item \ans{the sampling distribution of the sample mean}
          \item \ansline{Yes. Because the samples are random and come from a population with mean 100, the sample means should vary around 100 on average. The center of the sampling distribution equals the population mean.}
          \item \ansline{The spread (standard deviation) of the dot-plot would decrease. Larger samples reduce sampling variability because averaging more values smooths out the extremes.}
        \end{enumerate}

\end{enumerate}

\begin{extensionbox}
\textbf{Extension (optional).} \textit{Responses will vary. Look for: (1) identification of the two statistics and the claimed parameter, (2) acknowledgment that sampling variability alone can produce differences without a real population difference, and (3) some reasoning about sample size or context.}
\end{extensionbox}

\begin{reflectionbox}
\textbf{Preview:} Lesson 5.2 revisits the normal distribution, which will be the key
model for describing the shape of sampling distributions. Review your notes on
z-scores and normal probabilities from Unit 1.
\end{reflectionbox}

\end{document}
